% ==============================================================================
% KATA PENGANTAR (Pemerintah Desa Popontolen)
% ==============================================================================

\section*{Kata Pengantar}
\addcontentsline{toc}{section}{Kata Pengantar}

Puji syukur kami panjatkan ke hadirat Tuhan Yang Maha Esa, berkat rahmat dan karunia-Nya, publikasi \textbf{"Desa Popontolen Dalam Angka 2026"} dapat diselesaikan dengan baik. Publikasi ini disusun dalam rangka memperingati semangat kemerdekaan Republik Indonesia, sekaligus sebagai wujud nyata komitmen transparansi dan akuntabilitas data pembangunan desa berbasis potensi lokal.

Data yang disajikan dalam buku ini merupakan kompilasi data potensi wilayah dari 7 (tujuh) Jaga di Desa Popontolen, yang diharapkan dapat menjadi rujukan utama bagi Pemerintah Desa, Badan Permusyawaratan Desa (BPD), pemangku kepentingan (\textit{stakeholders}), serta seluruh elemen masyarakat dalam merumuskan kebijakan dan program pembangunan yang terarah, akurat, dan tepat sasaran.

Kami menyadari bahwa publikasi ini masih memiliki ruang untuk terus disempurnakan. Oleh karena itu, kritik, masukan, dan saran yang membangun sangat kami harapkan demi peningkatan kualitas dan keandalan data statistik desa di masa yang akan datang.

Apresiasi dan terima kasih yang setinggi-tingginya kami sampaikan kepada seluruh Kepala Jaga, perangkat desa, tim pembina statistik, dan masyarakat yang telah berpartisipasi aktif dalam proses pengumpulan dan penyusunan data ini. Semoga publikasi ini membawa manfaat sebesar-besarnya bagi kemajuan dan kesejahteraan masyarakat Desa Popontolen.

\vspace{1.5cm}
\begin{flushright}
  \begin{tabular}{c}
    Popontolen, Agustus 2026 \\
    \textbf{Hukum Tua / Kepala Desa Popontolen} \\[1.8cm]
    \textbf{\underline{Elke S. Poluakan, SKM, M.Kes}}
  \end{tabular}
\end{flushright}
